\scanpage{536}%
\mspage{447}%
% archive sheet number: 535
\[\left|\begin{aligned}&\operatorname{Ob}K=\operatorname{Ob}J\\ &\text{si }j,j'\in\operatorname{Ob}J,\ \text{on a } j\overset{K}{\leq}j'\quad\text{sss}\quad j\overset{J}{=}j',\ \text{ou }j\overset{J}{<}j'\ \text{et }\tau_{j',j}=p\end{aligned}\right.\]

Il est immédiat que c'est une \uline{relation d'ordre} sur $J$\inserted{(car $p^2=p$)}, plus fine que la relation d'ordre initiale.

\deleted{De plus, \deleted{\unsure{on|on a}} : \unsure{Car|car} si $i\overset{K}{<}j$ \deleted{\unsure{i.e.}} (i.e. $i<j$ et $\tau_{ji}=p$) alors}

\deleted{\[\left|\begin{aligned}j\leq j'&\Longrightarrow i\overset{K}{<}j'\\ i'\leq i&\Longrightarrow i'\overset{K}{<}j.\end{aligned}\right.\]}

\deleted{En effet,}

\deleted{\diagram{\begin{tikzcd}[ampersand replacement=\&] x_{\unsure{i|j}} \arrow[r] \arrow[d,equal] \& x_j\arrow[r]\arrow[d,equal] \& x_{j'}\\ x_0 \& \unsure{x|x_0} \& \end{tikzcd}}}

\deleted{on a}

\deleted{\[\begin{aligned}\tau_{j'i}&=\tau_{j'j}\tau_{ji}={}&\tau_{j'j}p\\ &&\Vert\ \text{par }1^\circ\\ &&p\tau_{j'j}\end{aligned}\]}

\deleted{\deleted{\unsure{On en a|On en a a}} donc $p\tau_{j'i}=\tau_{j'j}\unsure{p}$}

Le syst. inductif induit sur $K$ satisfait aux conditions $1^\circ,2^\circ,3^\circ$ de 4.15.7.3. \deleted{Il reste à} D'autre part $K\mathrel{\unsure{\longrightarrow|\relbar}}J$ est cofinal, donc aussi le composé $K\longrightarrow J\longrightarrow I$. Il reste à prouver que $K$ est grde devant $\pi_0$, donc que tte partie $K'\subset K$ telle que $\operatorname{card}K'\leq\pi_0$ admet un majorant dans $K$. Elle admet un majorant \uline{dans $\deleted{\unsure{I|J}}J$}, soit $j$. Par hypothèse 2 bis, il existe $j'\geq j$ tel que $\tau_{j'j}=p$ \deleted{\unsure{donc|et alors}}, et OPS $j'>j$. \deleted{\unsure{Sinon}} (En effet, si $\tau_{j'j}=p$, on aurait \deleted{que}
